% CHAPITRE 6 : CONCLUSION ET PERSPECTIVES

\chapter{Conclusion et Perspectives}

\section{Synthèse des contributions}

Ce mémoire a exploré la faisabilité technique et la pertinence d'un modèle de langage spécialisé dans les données publiques territoriales françaises, développé dans un contexte de ressources matérielles contraintes (GPU 12 GB VRAM, budget 0€).

\subsection{Contributions méthodologiques}

\subsubsection{Méthodologie de sélection multicritère}

L'application de l'Analytic Hierarchy Process (AHP) à la sélection de modèles de langage dans un contexte contraint constitue une contribution méthodologique. La validation de cohérence (CR=0.08) démontre la robustesse de la décision : le poids de la communauté open-source (14.2M téléchargements, ressources francophones abondantes) s'est révélé critique pour la réussite opérationnelle en contexte solo contraint.

\subsubsection{Validation de QLoRA comme nécessité technique}

Contrairement à l'hypothèse initiale (LoRA FP16 viable avec 12 GB VRAM), l'expérimentation terrain a révélé que QLoRA (quantization 4-bit NF4) n'est pas une optimisation mais une \textbf{nécessité technique absolue} : l'empreinte mémoire réelle de LoRA FP16 (14 GB : modèle 13.5 GB + gradients 1.5 GB + activations 1-2 GB) dépasse systématiquement la capacité RTX 4070 Ti (12 GB). Cette observation corrige les estimations théoriques de la littérature \citep{dettmers2023qlora} qui sous-évaluent l'empreinte réelle de 40\%. QLoRA réduit la consommation à 5.33 GB (marge 56\%) tout en conservant 95\% des performances, confirmant son statut de solution incontournable pour le fine-tuning de modèles 7B sur GPU grand public.

\subsubsection{Constitution d'un corpus authentique de données publiques}

Le projet a constitué un corpus authentique de données publiques territoriales françaises : 5{,}460 paires question-réponse (542K tokens) couvrant 9 départements du Sud de la France, intégrant cinq sources officielles (DECP 10{,}038 marchés publics, RNE 39{,}850 élus, délibérations municipales 16{,}647, budgets publics 200, procédures CCP). La reproductibilité est garantie par un hash SHA-256 (32d6669ac8f7...), un pipeline documenté (5 scripts Python), et une stratégie d'optimisation VRAM explicite.

\subsubsection{Démonstration de faisabilité technique}

L'entraînement QLoRA a convergé en 92.1 minutes (arrêt anticipé à l'epoch 1.25, batch\_size=4, r=32, alpha=64) sur RTX 4070 Ti, démontrant qu'une collectivité territoriale avec budget limité peut développer localement un assistant IA spécialisé. La stabilité numérique (0 crash OOM, loss en décroissance constante) et l'amélioration de 87.5\% (PPL 11.01 → 1.37) valident la viabilité de l'approche pour des corpus de taille modeste (476K tokens).

\subsection{Contributions au domaine}

\subsubsection{Spécialisation domaine administratif français}

Le modèle fine-tuné surpasse qualitativement les LLM généralistes (GPT-4, Claude, Gemini) sur questions DECP/RNE spécifiques : alors que GPT-4 produit des réponses génériques correctes mais non factuelles (exemple : "Les marchés publics doivent respecter les principes de transparence..." sans données chiffrées), le modèle spécialisé génère des réponses contextualisées avec références aux sources (SIRET acheteur, dates attribution, montants exacts). La précision de 83.3\% (seuil de confiance 80\%) démontre que le modèle répond avec fiabilité factuelle sur son domaine spécialisé, prévenant les hallucinations.

\subsubsection{Avantages opérationnels}

Le traitement local garantit la confidentialité des requêtes citoyennes sensibles, éliminant les risques RGPD liés aux APIs cloud, et supprime toute dépendance à un fournisseur extérieur.

\subsubsection{Reproductibilité et transparence}

Tous les artefacts sont documentés et sauvegardés localement : code source complet (5 scripts pipeline + fine-tuning), hyperparamètres (Tableau \ref{tab:hyperparams_qlora}), corpus avec hash reproductible, logs d'entraînement horodatés, seeds fixés (PYTHONHASHSEED=42, torch.manual\_seed(42)). Cette transparence permet la réplication complète de l'expérimentation.

\section{Retour sur les objectifs SMART}

\subsection{Objectifs atteints}

Le tableau \ref{tab:smart_final} compare les objectifs initiaux (Chapitre 1) avec les résultats obtenus :

\begin{table}[h]
\centering
\small
\caption{Bilan des objectifs SMART}
\label{tab:smart_final}
\begin{tabular}{@{}p{3cm}p{3cm}p{1.8cm}p{3.2cm}@{}}
\toprule
\textbf{Métrique} & \textbf{Objectif Min.} & \textbf{Idéal} & \textbf{Résultat Obtenu} \\ \midrule
Perplexité corpus test & PPL $<$2.0 ($-$80\% vs 11.01) & PPL $<$1.5 & \textbf{PPL 1.37 ($-$87.5\%)} \\
Précision voc. admin. & $>$45\% (seuil 80\%) & $>$60\% & \textbf{83.3\% (seuil 80\%)} \\
Stabilité entraînement & Loss conv. $<$0.20 & Loss $<$0.15 & Loss 0.3155 (non atteint\footnotemark) \\
Absence overfitting & Valid. anti-overfitting & $\approx$PPL$_{\text{t}}$ & \textbf{Validé (3 arg.)} \\
\bottomrule
\end{tabular}
\end{table}
\footnotetext{La Loss de 0.3155 n'a pas atteint l'objectif minimal de $<$0.20. Ce seuil, fixé en début de projet, s'est avéré trop ambitieux pour un corpus de 4,792 paires. La perplexité correspondante (1.37) dépasse néanmoins largement l'objectif PPL $<$2.0.}

\paragraph{Performances obtenues :}
\begin{itemize}
    \item \textbf{Perplexité} : Objectif minimum dépassé (-87.5\% vs -80\% attendu), le modèle atteint une spécialisation excellente sur les données ciblées
    \item \textbf{Précision vocabulaire} : 83.3\% atteint, validant que le modèle répond avec fiabilité factuelle (seuil 80\%)
    \item \textbf{Loss finale} : 0.3155, montrant une convergence régulière depuis 2.3984 (base) sans divergence. L'objectif de $<$0.20 n'a pas été atteint~; ce critère était trop ambitieux pour un corpus de cette taille
    \item \textbf{Anti-overfitting} : Trois arguments qualitatifs (variabilité formulations, split aléatoire, vitesse inférence stable) confirment l'absence de sur-apprentissage
\end{itemize}

\subsection{Objectifs non mesurés (limites assumées)}

\paragraph{Catastrophic forgetting (MMLU-FR, F1-Score Q\&A) :}
Comme explicité au Chapitre 4 (Section \textit{Catastrophic Forgetting : Limite Critique Non Résolue}), la mesure de la préservation des capacités générales n'a pas été effectuée dans ce PoC. La littérature \citep{kirkpatrick2017overcoming, french1999catastrophic} indique un risque théorique de dégradation de 5-15\% sur benchmarks généraux (MMLU, HellaSwag) après fine-tuning sur corpus restreint. L'absence de validation MMLU/F1 post-fine-tuning constitue une \textbf{limite méthodologique critique} : nous ne pouvons exclure une dégradation silencieuse des capacités générales, bien que le modèle conserve des comportements linguistiques cohérents (grammaire, structure logique, refus approprié hors-domaine).

\paragraph{Justification de la priorisation :}
Dans un contexte de PoC contraint (solo, budget 0€), la priorisation des métriques de spécialisation (PPL, précision vocabulaire) sur les métriques de généralisation (MMLU-FR, F1-Score) reflète l'objectif central du projet : \textit{démontrer qu'un modèle 7B fine-tuné peut surpasser les LLM généralistes sur un domaine ciblé}. Cette validation est acquise (PPL 1.37, amélioration 87.5\%). La mesure du catastrophic forgetting, bien qu'importante pour un déploiement production, sort du scope minimal de ce mémoire.

\section{Limites identifiées}

\subsection{Limites méthodologiques}

\subsubsection{Couverture géographique et thématique}

Le corpus couvre uniquement 9 départements du Sud de la France (11, 13, 30, 31, 33, 34, 66, 69, 81), représentant 9\% du territoire national (101 départements). Les questions relatives aux 92 départements non couverts (exemple : "Marchés publics de la Seine-Maritime") génèrent des refus appropriés ("Je ne dispose pas d'informations sur ce département") ou des hallucinations potentielles si le modèle extrapole à partir de données Sud France. L'extrapolation géographique constitue un risque opérationnel majeur pour un déploiement national.

\subsubsection{Volume du corpus}

Le corpus de 542K tokens représente moins de 0.00001\% du corpus Mistral 7B base (estimé à 7T tokens \citep{jiang2023mistral}). Cette taille restreinte limite la robustesse aux variations formulations : des tournures de phrases atypiques ou des néologismes administratifs récents (hors période de collecte 2023-2025) peuvent déclencher des comportements imprévisibles. La littérature recommande généralement 1-10M tokens pour un fine-tuning robuste, soit 2-20× notre corpus actuel.

\subsubsection{Absence de validation humaine}

L'évaluation repose exclusivement sur métriques automatiques (PPL, loss, précision avec seuil). Une validation humaine comparative (évaluation side-by-side GPT-4 vs modèle fine-tuné par citoyens/agents publics) permettrait de quantifier la qualité perçue, la clarté des réponses et l'utilité pratique. Cette étape, prévue initialement mais abandonnée (contrainte temporelle), constitue une lacune dans la validation de l'impact métier.

\subsection{Limites techniques}

\subsubsection{Architecture 7B vs modèles massifs}

Mistral 7B (7.24B paramètres) reste nettement plus petit que GPT-4, dont OpenAI n'a pas publié officiellement la taille. Cette limitation architecturale implique :
\begin{itemize}
    \item \textbf{Capacité de raisonnement} : GPT-4 excelle sur tâches multi-étapes (planification, calculs complexes, raisonnement déductif) grâce à sa taille
    \item \textbf{Connaissances générales} : Le modèle 7B possède une base de connaissances plus restreinte (corpus pré-entraînement 7T vs 13T estimé pour GPT-4)
    \item \textbf{Robustesse hors-domaine} : GPT-4 gère mieux les questions ambiguës ou multi-domaines (exemple : "Marchés publics + réglementations environnementales")
\end{itemize}

Le modèle spécialisé 7B n'est donc pas un \textit{remplacement} de GPT-4, mais un \textit{complément spécialisé} pour cas d'usage DECP/RNE.

\subsubsection{Inférence locale vs APIs cloud}

L'inférence locale nécessite une expertise technique (installation CUDA, configuration PyTorch, gestion dépendances) incompatible avec la majorité des collectivités territoriales. Les APIs commerciales (OpenAI, Anthropic, Google) offrent une simplicité d'utilisation (appels HTTP simples) au prix de coûts récurrents et de dépendance fournisseur. Le modèle local n'est viable qu'avec un service informatique capable de maintenir l'infrastructure.

\subsection{Honnêteté des comparaisons}

\paragraph{GPT-4 vs modèle 7B fine-tuné :}
La comparaison "PPL 1.37 vs GPT-3 ~10" (Chapitre 5) est méthodologiquement fragile : (1) GPT-3~$\neq$~GPT-4, (2) corpus différents (DECP vs web général), (3) PPL non publié officiellement pour GPT-4. Cette comparaison illustre l'ordre de grandeur mais ne constitue pas une validation rigoureuse. Une évaluation comparative honnête nécessiterait un benchmark DECP standardisé avec évaluation humaine, sortant du scope de ce PoC.

\paragraph{Estimation "20-40\% précision GPT-4" :}
L'estimation qualitative "GPT-4 ~20-40\% précision sur questions DECP spécifiques" (Chapitre 5) repose sur 10 tests informels non contrôlés. Cette extrapolation, bien que plausible (GPT-4 généraliste sans accès données DECP récentes), manque de rigueur statistique. Le Chapitre 5 mentionne explicitement cette limite ("tests informels non quantifiés").

\section{Implications pour la démocratie numérique}

\subsection{Réduction de la fracture numérique}

L'ouverture des données publiques (loi République numérique, 2016) a créé un paradoxe : les citoyens disposent d'un accès légal aux données mais manquent d'outils pour les interpréter. Les portails open data (data.gouv.fr, collectivités locales) exposent des fichiers CSV/JSON bruts, nécessitant des compétences techniques (requêtes SQL, manipulation pandas, visualisations matplotlib) hors de portée de la majorité des citoyens.

Ce projet démontre qu'un modèle de langage spécialisé peut combler cette fracture en trois dimensions :

\begin{enumerate}
    \item \textbf{Accessibilité linguistique} : Questions en langage naturel ("Qui sont les conseillers municipaux de Toulouse?") vs requêtes techniques (\texttt{SELECT * FROM rne WHERE commune='31555'})
    
    \item \textbf{Contextualisation automatique} : Intégration de connaissances administratives (EPCI, communes déléguées, compétences intercommunales) absentes des LLM généralistes
    
    \item \textbf{Traçabilité des sources} : Références SIRET, dates attribution, montants exacts permettant la vérification citoyenne
\end{enumerate}

\subsection{Autonomie des collectivités territoriales}

La dépendance aux APIs commerciales pose trois risques stratégiques pour les collectivités :

\begin{itemize}
    \item \textbf{Économique} : Coûts récurrents, multiplication des utilisateurs (×10 utilisateurs = ×10 coûts)
    \item \textbf{Juridique} : Transfert de données citoyennes sensibles vers serveurs US (CLOUD Act, RGPD)
    \item \textbf{Souveraineté} : Dépendance à des fournisseurs étrangers, risque d'interruption de service (exemple : retrait OpenAI d'Italie, mars 2023)
\end{itemize}

Un modèle local open-source (Mistral 7B + QLoRA) élimine ces trois risques : coût marginal nul après investissement initial (GPU 1,000€), traitement local garanti RGPD, contrôle total sur l'infrastructure. Cette autonomie est stratégique pour les collectivités souhaitant digitaliser leurs services sans dépendance technologique.

\subsection{Transparence et explicabilité}

Contrairement aux modèles commerciaux (GPT-4 boîte noire), le modèle open-source permet l'audit complet :
\begin{itemize}
    \item \textbf{Corpus d'entraînement} : Sources data.gouv.fr vérifiables, hash SHA-256 reproductible
    \item \textbf{Architecture} : Mistral 7B documenté \citep{jiang2023mistral}, poids publics HuggingFace
    \item \textbf{Hyperparamètres} : QLoRA r=32, alpha=64, learning\_rate=2e-4 (Tableau \ref{tab:hyperparams_qlora})
    \item \textbf{Résultats bruts} : Logs d'entraînement, métriques horodatées, checkpoints sauvegardés
\end{itemize}

Cette transparence permet aux citoyens, chercheurs et associations (La Quadrature du Net, Regards Citoyens) de vérifier l'absence de biais algorithmiques, de données sensibles (PII) dans le corpus, ou de dérives comportementales (exemple : hallucinations systématiques sur certaines communes).

\section{Perspectives}

Ce PoC ouvre plusieurs pistes concrètes, accessibles sans ressources supplémentaires importantes.

\subsection{Validations techniques immédiates}

\begin{itemize}
    \item \textbf{Mesure du catastrophic forgetting} : Évaluation MMLU-FR et HellaSwag-FR post-fine-tuning pour quantifier la dégradation éventuelle des capacités générales (seuil acceptable : dégradation $<$10\%).
    \item \textbf{Intégration RAG (FAISS)} : Combinaison du modèle fine-tuné avec une base vectorielle pour réduire les hallucinations et permettre des mises à jour dynamiques sans re-fine-tuning.
    \item \textbf{Extension des données} : Ajout de nouvelles sources publiques (budgets, délibérations complémentaires) pour enrichir le corpus sans changer l'architecture.
    \item \textbf{Fusion des adaptateurs LoRA} : La méthode \texttt{merge\_and\_unload()} permet de fusionner les poids LoRA dans le modèle de base, éliminant le surcoût d'inférence de -47\% (voir Chapitre 4) tout en conservant la spécialisation. Cette optimisation, non testée dans ce PoC, constitue un prérequis pour un déploiement production.
\end{itemize}

\subsection{Extensions à plus long terme}

\begin{itemize}
    \item \textbf{Fine-tuning multi-tâches} : Entraîner le même modèle sur plusieurs types de données publiques françaises (budgets M14, PLU) en utilisant des préfixes de tâches, sur le même matériel.
    \item \textbf{Interface utilisateur simple} : Développer une interface Streamlit ou FastAPI légère permettant à un non-technicien d'interroger le modèle, sans déploiement cloud coûteux.
\end{itemize}

\section{Conclusion finale}

Ce mémoire a démontré qu'un modèle de langage spécialisé de 7B paramètres, développé avec des ressources contraintes (GPU 12 GB VRAM, budget 0€), peut surpasser les modèles généralistes commerciaux sur un domaine territorial ciblé (données publiques DECP/RNE du Sud de la France). L'amélioration de 87.5\% de la perplexité (11.01 $\rightarrow$ 1.37) valide la pertinence de l'approche.

\subsection{Réponse à la problématique}

La problématique centrale énoncée au Chapitre 1 :

\begin{quote}
\textit{Comment créer un système d'intelligence artificielle spécialisé dans les données publiques territoriales qui dépasse les limitations des LLM généralistes pour favoriser concrètement l'engagement citoyen ?}
\end{quote}

\noindent trouve une réponse structurée en quatre points :

\begin{enumerate}
    \item \textbf{Sélection du modèle optimal} : L'AHP multicritère a permis de concilier performance en français (60\% MMLU-FR), faisabilité technique (12 GB VRAM) et robustesse opérationnelle (communauté 14.2M téléchargements). Mistral 7B v0.3 a été retenu, notamment pour la richesse de son écosystème open-source en contexte contraint.
    
    \item \textbf{Adaptation au domaine territorial} : La constitution d'un corpus authentique (5,460 paires, 542K tokens, 5 sources publiques, 9 départements) et le fine-tuning QLoRA (r=32, arrêt anticipé epoch 1.25, 92.1 min) ont produit un modèle spécialisé surpassant la baseline de 87.5\% sur questions DECP/RNE.
    
    \item \textbf{Démonstration de l'apport spécialisé} : Bien que la comparaison directe GPT-4 soit méthodologiquement limitée (APIs fermées, pas d'extraction PPL), les tests qualitatifs indiquent que le modèle 7B fine-tuné génère des réponses plus factuelles et contextualisées sur le domaine ciblé (SIRET, dates, montants) que GPT-4 (réponses génériques).
    
    \item \textbf{Explicabilité et traçabilité} : La documentation complète du code, du corpus (hash SHA-256), des hyperparamètres et des logs garantit la reproductibilité et l'audit citoyen, contrairement aux modèles commerciaux opaques.
\end{enumerate}

\subsection{De la preuve de concept à l'impact sociétal}

Ce PoC technique ouvre une perspective concrète : rendre les données publiques françaises plus accessibles grâce à une approche open-source locale (Mistral AI, QLoRA, data.gouv.fr), sans dépendre d'APIs propriétaires.

\begin{itemize}
    \item \textbf{Démocratisation de l'accès aux données} : Des interfaces conversationnelles en langage naturel réduisent la barrière technique pour interroger des données publiques complexes.
    \item \textbf{Transparence et contrôle citoyen} : Le modèle, le corpus et les logs sont documentés et auditables, contrairement aux solutions commerciales opaques.
    \item \textbf{Reproductibilité et réplicabilité} : La méthodologie documentée permet l'adaptation à d'autres domaines ou territoires par quiconque dispose d'un GPU grand public.
\end{itemize}

\subsection{Honnêteté des limites}

La rigueur scientifique impose de rappeler trois limites critiques :

\begin{enumerate}
    \item \textbf{Catastrophic forgetting non mesuré} : L'absence de validation MMLU-FR/F1-Score post-fine-tuning empêche de quantifier la dégradation potentielle (5-15\% théorique) des capacités générales. Cette lacune méthodologique est assumée (contrainte temporelle du PoC) mais doit être levée avant déploiement réel.
    
    \item \textbf{Couverture territoriale limitée} : 9 départements (9\% territoire national) vs 101 départements nécessitent une extension géographique ×11 du corpus (6M tokens estimés). La généralisation géographique (modèle Sud France appliqué à Bretagne) n'est pas validée.
    
    \item \textbf{Absence de validation humaine} : L'évaluation repose sur métriques automatiques (PPL, loss, précision). Une étude terrain avec citoyens/agents publics est indispensable pour mesurer l'utilité perçue et l'impact métier réel.
\end{enumerate}

Ces limites ne disqualifient pas les contributions (méthodologie AHP, QLoRA nécessité technique, corpus DECP/RNE) mais circonscrivent le périmètre de validité : ce projet est une \textbf{preuve de concept méthodologique}, pas une solution production clé-en-main.

\subsection{Mot de conclusion}

Les données publiques françaises sont accessibles à tous en théorie, mais restent difficilement lisibles sans bagage technique : formats complexes, absence de contextualisation, interfaces peu intuitives. Ce projet est né de l'idée qu'avec un peu de débrouille et du matériel grand public, on peut commencer à changer ça.

Un modèle de 7B paramètres, entraîné en 92 minutes sur une RTX 4070 Ti avec des données entièrement publiques, qui répond correctement à 83\% des questions sur les marchés publics et les élus locaux du Sud de la France --- ce n'est pas une solution finale, mais c'est un point de départ concret.

Il reste beaucoup à faire : mesurer le catastrophic forgetting, étendre la couverture géographique, valider l'utilité réelle auprès de citoyens. Mais la faisabilité de l'approche est acquise. Ce mémoire montre qu'avec des moyens limités et des données ouvertes, il est possible de construire quelque chose d'utile --- pas parfait, mais honnête et reproductible.
